% Apêndice — protocolo da busca bibliográfica do Capítulo de Trabalhos
% Relacionados (task T02 da revisão v1). Strings, contagens e critérios
% transcritos do registro que estava em comentário em
% 2-textuais/trabalhos-relacionados.tex (busca de 02/11/2026).
\apendice{Protocolo da busca bibliográfica}
\label{ap:busca}

Este apêndice registra a busca que sustenta o Capítulo~\ref{cap:trabalhos-relacionados}
e a afirmação, da Seção~\ref{sec:problema}, de que não foram encontrados
estudos que reúnam as características da proposta. A busca foi feita em
2 de novembro de 2026, com filtro de período de 2023 a 2026, exceto para os
dois trabalhos fora do domínio da IA (\citeonline{kaur2017};
\citeonline{feitosa2017}), escolhidos pelo tema.

\section*{Bases com \textit{strings} registradas}

Na ACM Digital Library, acessada pelo Portal de Periódicos CAPES (acesso UFC),
as \textit{strings} foram aplicadas ao campo \textit{Abstract}, exceto a A3,
aplicada ao título. Os resultados são dados para a \textit{ACM Full-Text
Collection} (\textit{dl}) e, quando registrado, para o \textit{ACM Guide}
(\textit{guide}), que inclui outras editoras.

\begin{description}
  \item[A1] (``large language model'' OR LLM) AND (``vendor lock-in'' OR
  ``abstraction layer'' OR middleware OR gateway OR ``provider-agnostic'' OR
  ``model-agnostic'') AND (architecture OR integration). \textit{dl}: 61.
  \item[A2] (``large language model'' OR LLM OR ``language model'') AND
  (``vendor lock-in'' OR ``lock-in'' OR ``provider-agnostic'' OR ``switching
  provider'' OR ``provider switching'' OR ``model migration'' OR portability OR
  replaceability OR substitutability). \textit{dl}: 439.
  \item[A3] Título: (LLM OR ``language model(s)'' OR ``foundation model'') AND
  Título: (gateway OR middleware OR abstraction OR ``lock-in'' OR portability OR
  ``provider-agnostic'' OR migration). \textit{dl}: 82; \textit{guide}: 194.
  \item[A4] (``large language model'' OR LLM) AND (``software architecture'' OR
  ``architectural decision'' OR ``reference architecture'' OR ``architectural
  tactic'') AND (integrated OR integration). \textit{dl}: 10; \textit{guide}: 20.
  \item[A5] energy AND (``agentic framework(s)'' OR ``agent framework'' OR
  LangChain OR AutoGen OR ``orchestration framework'' OR ``LLM framework'').
  \textit{dl}: 7; \textit{guide}: 26.
  \item[A6] energy AND (``design pattern(s)'' OR ``architectural tactic(s)'' OR
  ``abstraction layer'' OR indirection) AND (consumption OR overhead OR
  efficiency). \textit{dl}: 33; \textit{guide}: 118.
  \item[A7] (``large language model'' OR LLM) AND (LangChain OR LiteLLM OR ``LLM
  framework(s)'' OR ``LLM library'' OR ``LLM SDK'' OR ``LLM API'') AND (empirical
  OR developers OR maintenance OR evolution). \textit{dl}: 52; \textit{guide}: 103.
\end{description}

No IEEE Xplore, também pelo Portal CAPES, no campo \textit{Abstract}:

\begin{description}
  \item[I1] a mesma \textit{string} da A1. Resultados: 90.
  \item[I2] (energy OR ``power consumption'') AND (``large language model'' OR
  LLM OR agentic) AND (framework OR orchestration OR middleware OR
  ``abstraction layer'' OR ``software architecture'') AND (overhead OR ``energy
  consumption''). Resultados: 171.
\end{description}

\section*{Buscas complementares}

A OpenAlex e a Crossref (por suas APIs) e a busca na \textit{web} serviram
para localizar trabalhos por tema e conferir metadados. A SBC OpenLib foi
buscada por evento, e não pela busca geral, que exige \textit{login}: SBES,
SBCARS, SEMISH, WSCAD/SSCAD, WCAMA e SBSI. Para essas buscas não há
\textit{strings} nem contagens registradas, de modo que a repetição exata do
protocolo vale para a ACM Digital Library e para o IEEE Xplore.

\section*{Triagem}

A triagem foi feita por título e resumo. O critério de inclusão foi o trabalho
ser revisado por pares e ter a mesma motivação, o mesmo objeto ou o mesmo
método da proposta. Foram excluídos os \textit{preprints} sem veículo de
publicação, as revisões redundantes com um trabalho já incluído, os trabalhos
que comparam \textit{frameworks} apenas pela qualidade da saída do modelo e os
resultados fora do tema (aplicações de LLMs, segurança, inferência em
\textit{hardware}). Dos resultados triados, 21 trabalhos foram incluídos e
constam do Quadro~\ref{qua:trabalhos-relacionados}; o número de trabalhos
triados e excluídos em cada etapa não foi registrado.
